% TOP_PROBLEM: {"id": 3152, "title": "Nearly spanning regular subgraphs", "category_id": 3, "category": {"id": 3, "name": "graph_theory", "display_name": "Graph Theory", "description": "Problems involving graphs, networks, and their properties.", "slug": "graph-theory", "order_index": 3, "created_at": "2026-07-31T15:26:25.671Z"}, "status": "open", "problem_number": "OPG-801", "set_id": 12}
% FIRST_POSTED: 2026-10-01
% ENTRY_KIND: statement_only
% UnsolvedMath ID 3152; source catalogue code OPG-801
% !TeX program = lualatex
% Definition and sources only; this file is not an LLM solution attempt.
\documentclass[11pt,a4paper]{article}
\usepackage[margin=25mm]{geometry}
\usepackage{fontspec}
\setmainfont{lmroman10-regular.otf}[BoldFont=lmroman10-bold.otf,
 ItalicFont=lmroman10-italic.otf,BoldItalicFont=lmroman10-bolditalic.otf]
\usepackage{amsmath,amssymb,mathtools}
\usepackage{unicode-math}
\setmathfont{latinmodern-math.otf}
\AtBeginDocument{\let\setminus\smallsetminus}
\usepackage{xurl,hyperref}
\hypersetup{colorlinks=true,urlcolor=blue}
\setcounter{secnumdepth}{0}
\setlength{\parindent}{0pt}
\setlength{\parskip}{5pt}
\setlength{\emergencystretch}{3em}
\begin{document}
\section*{Nearly spanning regular subgraphs}\label{3152}
{\small UnsolvedMath ID 3152 · OPG-801 · Graph Theory · OpenGarden}
\subsection{Definitions and mathematical statement}
A finite simple graph $G$ is $r$-regular if every vertex has exactly
$r$ neighbours. A $k$-regular subgraph $H$ consists of a vertex set
$U\subseteq V(G)$ and an edge set $F\subseteq E(G[U])$ such that
every vertex of $(U,F)$ has degree exactly $k$. Edges inside $U$
may be omitted, so $H$ need not be induced. It also need not be
connected.

The conjecture states that for every integer $k\geq1$ and every
real $\varepsilon$ with $0<\varepsilon<1$, there is an integer
$r_0=r_0(k,\varepsilon)\geq k$ such that, for every integer
$r\geq r_0$ and every nonempty finite simple $r$-regular graph $G$,
one can find a $k$-regular subgraph $H$ satisfying
\[
 |V(H)|\geq(1-\varepsilon)|V(G)|.
\]
Restricting to $\varepsilon<1$ gives the substantive part of the
source's all-positive-$\varepsilon$ formulation; larger errors
follow from any smaller positive error.

The threshold may depend on the desired degree and the fraction
of vertices one is prepared to discard, but not on the order,
connectedness, or other features of $G$. The conclusion demands
exact regularity on the retained vertices, not merely average
degree $k$ or degrees close to $k$. For odd $k$, the handshaking
identity forces $|V(H)|$ to be even; the permitted vertex deletion
allows this necessary parity condition to be met.
\subsection{Short English statement}
Does every sufficiently high-degree regular simple graph contain an exactly prescribed-degree regular subgraph covering almost all vertices?
\subsection{Sources}
\begin{itemize}
\item[S1] ULAM.AI and the UnsolvedMath Contributors, supplied problems.json, record 3152 (OPG-801), OpenGarden collection. Catalogue identity and source transcription; CC BY 4.0.
\url{https://huggingface.co/datasets/ulamai/UnsolvedMath}.
\item[S2] Open Problem Garden, Nearly spanning regular subgraphs. Original problem and discussion; the mathematical formulation below is expanded from the supplied source record.
\url{http://www.openproblemgarden.org/op/nearly_spanning_regular_subgraphs}.
\end{itemize}
\end{document}
